%%%%%%%%%%%%%%%%%%%%%%%%
% $Autor: RanjeetsingTawar$
% $Pfad: Car_Licence_Plate_Identification_with_a_JetsonNano$
% $Version: 4.0.0 $
%
% !TeX encoding = utf8
%
%%%%%%%%%%%%%%%%%%%%%%%%

\chapter{Pandas package}

\section{Introduction}

Pandas is a widely-used, open-source Python library designed for data analysis and manipulation. It offers user-friendly data structures and tools that make working with structured data, such as spreadsheets, SQL databases, and time series, highly efficient.

This report provides a detailed overview of the Pandas library, covering its features, installation process, and practical examples that illustrate its application in various scenarios. Additionally, resources for further exploration of the library will be recommended for readers who wish to deepen their understanding \cite{McKinney:2011}. 

The version of Pandas discussed in this report is 1.4.4. Within the context of the project, Pandas is utilized to load data from CSV files, manage missing values, create new columns, and perform other transformations. These functionalities are explained in detail in the sections that follow.\cite{McKinney:2015}.

\section{Description}

Pandas is a highly popular and versatile open-source Python library for data analysis and manipulation. It offers intuitive data structures and tools, making it exceptionally efficient for handling structured data such as spreadsheets, SQL databases, and time series.

The foundation of Pandas lies in its two main data structures: Series and DataFrame.

\begin{itemize}
	
	\item A Series in Pandas is a one-dimensional labeled array capable of storing data of any type. It functions similarly to a column in a spreadsheet or a database table and allows data to be accessed through labels or indices. The Series structure enables efficient data manipulation and analysis, supporting operations like filtering, aggregation, and transformation \cite{McKinney:2011}.
	
	\item A DataFrame, in contrast, is a two-dimensional labeled data structure with columns that can hold different data types. It resembles a tabular format, akin to a spreadsheet or an SQL table. DataFrames offer a comprehensive range of operations, including indexing, merging, reshaping, and grouping, enabling powerful data manipulation and analysis.
	
\end{itemize}

Pandas provides an extensive array of features, such as data cleaning, transformation, filtering, merging, reshaping, and visualization. It includes tools for managing missing data, processing time series, and conducting statistical operations. With its user-friendly and expressive syntax, Pandas streamlines data manipulation tasks and enhances the efficiency of data analysis workflows \cite{Heydt:2017}.

Additionally, Pandas integrates effortlessly with other widely-used Python libraries like NumPy, Matplotlib, and scikit-learn, establishing itself as a cornerstone of the data science and analytics ecosystem. Its flexibility, performance, and comprehensive documentation make it a top choice for data professionals, researchers, and analysts working with structured data in Python.

\subsection{Data Suitability}

Pandas is highly compatible with various types of data.

\begin{itemize}
	
	\item \textbf{Data Structure:} Pandas operates primarily with two core data structures: Series and DataFrame. A Series is a one-dimensional labeled array that can store data of any type, whereas a DataFrame is a two-dimensional labeled structure similar to a table or spreadsheet. To use Pandas effectively, ensure your data is organized in a format compatible with these structures.
	
	\item \textbf{Data Format:} Pandas supports a wide range of data formats, such as CSV, Excel, SQL databases, JSON, HTML, HDF5 (Hierarchical Data Format), Parquet, Feather, and others. To use Pandas effectively, ensure your data is in a compatible format that can be easily imported. Additionally, Pandas offers functions to read and write data across various formats, enabling seamless interaction with diverse data sources.
	
	\item \textbf{Data Size:} When working with Pandas, it’s important to consider the size of your dataset. While Pandas is capable of handling large datasets, extremely large ones may lead to performance bottlenecks or memory limitations. For such cases, you may need to optimize your code, select efficient data types, or explore alternative tools like \textbf{Dask} or \textbf{Apache Spark} for distributed computing.
	
	\item \textbf{Performance Considerations:} The size and complexity of your data can impact the computational cost of certain operations in Pandas. To enhance performance, it is crucial to optimize your code by utilizing Pandas' built-in vectorized operations and efficient indexing methods.
	
	
\end{itemize}

\subsection{Key Features of Pandas}

Pandas excels in the following areas \cite{Betancourt:2019}:

\begin{itemize}
	
	\item Efficient handling of missing data (represented as NaN) in both floating-point and non-floating-point data types.
	
	\item Size mutability, enabling the addition and removal of columns from DataFrames and higher-dimensional objects.
		
	\item Automatic and explicit data alignment, which allows manual alignment of objects to a set of labels or automatic alignment during calculations involving Series, DataFrames, etc.
	
	\item Powerful and flexible group by functionality, supporting split-apply-combine operations for both data aggregation and transformation.
	
	\item Seamless conversion of ragged and differently-indexed data from other Python and NumPy data structures into DataFrame objects.
	
	\item Intelligent label-based slicing, fancy indexing, and subsetting of large data sets.
	
	\item Intuitive merging and joining of data sets.
	
	\item Flexible reshaping and pivoting capabilities for datasets.
	
	\item Hierarchical labeling of axes, allowing the presence of multiple labels per tick.
	
	\item Robust I/O functions for importing data from flat files (CSV and delimited formats), Excel files, databases, and efficient reading and writing of data in the HDF5 format.
	
	\item Time series-specific functionality, such as generating date ranges, converting frequencies, calculating moving window statistics, date shifting, and lagging.
	
\end{itemize}

\subsection{Why Multiple Data Structures in Pandas?}


Pandas uses various data structures, each serving a particular purpose. It is helpful to think of Pandas data structures as flexible containers for lower-dimensional data. For example, a DataFrame serves as a container for Series, while a Series acts as a container for scalars. This design allows for the insertion and removal of objects from these containers in a way similar to dictionaries \cite{McKinney:2015}.

Pandas also offers intuitive default behaviors for common API functions, taking into account the typical structure of time series and cross-sectional data. When using N-dimensional arrays (ndarrays) for storing 2- and 3-dimensional data, the user must manage the orientation of the dataset when writing functions. In contrast, Pandas' axis design adds semantic meaning to the data, making it easier to understand. For any given dataset, there is usually a preferred orientation, aiming to reduce the cognitive load required when performing data transformations in subsequent functions \cite{McKinney:2012}.

For example, when working with tabular data (DataFrame), it is more meaningful to think in terms of the index (representing rows) and columns rather than axis 0 and axis 1. Iterating through the columns of a DataFrame leads to clearer and more readable code.

\begin{code}[h!]
	\lstinputlisting[language=Python, numbers=none, linerange={30-32}]{Code/Pandas/Pandas.py}    
	
	\caption{Columns of DataFrame}
	
\end{code}

\section{Installation}

The most efficient way to install Pandas is by using the cross-platform Anaconda distribution, which is specifically designed for data analysis and scientific computing. This installation method is highly recommended for the majority of users due to its convenience and ease of use.

\subsection{Python version support}

Pandas officially supports a range of Python versions, including Python 3.8, 3.9, 3.10, and 3.11.

\subsection{Installing with Anaconda}

For users with less experience, installing Pandas along with the entire NumPy and SciPy stack can be difficult. However, Anaconda, a cross-platform (Linux, macOS, Windows) Python distribution designed for data analytics and scientific computing, simplifies the process. With Anaconda, users can easily install not just Pandas, but also Python and other essential tools in the SciPy stack (such as IPython, NumPy, Matplotlib, etc.). Once the installer is launched, there is no need for additional installations or waiting for software compilation to begin using Pandas and the rest of the SciPy stack.


\subsection{Installing using terminal or command prompt}

To install Pandas, you can use the Python package manager, pip, by running the following command in your terminal or command prompt:

\begin{code}[h!]
	\lstinputlisting[language=Python, numbers=none, linerange={34-35}]{Code/Pandas/Pandas.py}    
	
	\caption{Installing Pandas using terminal}
	
\end{code}	

Ensure that you have a compatible version of Python installed on your system before running the installation command.

\subsection{Required dependencies}

Pandas requires the following dependencies:

\begin{center}
	\begin{tabular}{ |c|c| }
		\hline
		Package & Minimum supported version \\
		\hline
		NumPy & 1.20.3 \\
		python-dateutil & 2.8.2 \\
		pytz & 2020.1 \\
		\hline
	\end{tabular}
\end{center}

\section{Example - Manual}

\subsection{User Manual for the Example Python File in PyCharm}

\subsubsection{Installation and Setup}

\begin{enumerate}
	
	\item \textbf{Download and install Python:}
	
	Visit the official Python website \url{https://www.python.org} and download the latest version of Python for your operating system. Follow the provided installation instructions.
	
	\item \textbf{Install PyCharm: }
	
	Visit the JetBrains website \url{https://www.jetbrains.com/pycharm} and download the free PyCharm Community Edition. Follow the installation steps according to your operating system.
	
\end{enumerate}

\subsubsection{Opening the Python File in PyCharm}

\begin{enumerate}
	
	\item Launch PyCharm: Open the PyCharm application from your desktop or applications menu.
	
	\item Create a new project: Click on \texttt{"Create New Project"} or go to \texttt{File} $\rightarrow$ \texttt{"New Project"}. Choose a suitable location for your project and provide a name.
	
	\item Open the Python file: Once the project is created, navigate to the project directory in the PyCharm project view. Right-click on the desired folder and select \texttt{"New"} $\rightarrow$ \texttt{"Python File"}. Provide a name for the file and click \texttt{"OK"}.
	
	\item Copy your Python code: Open your Python file (with .py extension) in a text editor and copy the contents.
	
	\item Paste the code: Paste the copied code into the newly created Python file in PyCharm.
	
\end{enumerate}

\subsubsection{Working with the Python File}

\begin{enumerate}
	
	\item Running the script: To run the Python script, right-click anywhere within the Python file and select \texttt{"Run"} $\rightarrow$ \texttt{"Run ExampleManual.py"}. Alternatively, you can use the keyboard shortcut \texttt{"Shift + F10"}. The script will execute, and the output will be displayed in the PyCharm console.
	
	\item Debugging the script: To debug the Python script and set breakpoints for analysis, click on the left gutter of the Python file, next to the line where you want to set the breakpoint. A red dot will appear, indicating the breakpoint. Click on the \texttt{"Debug"} button or use the keyboard shortcut \texttt{"Shift + F9"} to start debugging the script.
	
	\item Interacting with the script: If your script expects user input or provides interactive prompts, you can provide the input in the PyCharm console. The console allows you to interact with the script while it is running.
	
\end{enumerate}

\subsubsection{Modifying the Python File}

\begin{enumerate}
	
	\item Editing the code: To make changes to the Python code, simply locate the section you want to modify and edit the code accordingly.
	
	\item Saving the changes: PyCharm automatically saves your changes as you work. However, you can manually save the file by going to \texttt{"File"} $\rightarrow$ \texttt{"Save"} or using the keyboard shortcut \texttt{"Ctrl + S"}.
	
\end{enumerate}

\subsubsection{Further Assistance and Resources}

\begin{itemize}
	
	\item PyCharm Documentation: PyCharm offers comprehensive documentation to help you understand its features and functionality. You can access it online at \texttt{https://www.jetbrains.com/help/pycharm}.
	
	\item Python Documentation: The official Python documentation provides detailed information about the Python language, libraries, and best practices. It is available at \url{https://docs.python.org}.
	
	\item Online Python Communities: Joining online communities like Stack Overflow \url{https://stackoverflow.com} or the Python subreddit \url{https://www.reddit.com/r/Python} can provide valuable insights and assistance from experienced Python developers.
	
\end{itemize}


\subsection{How to import}
The code begins by importing the required libraries, numpy (\texttt{np}) and pandas (\texttt{pd}).

\begin{code}[h!]
	\lstinputlisting[language=Python, numbers=none, linerange={37-38}]{Code/Pandas/Pandas.py}    
	
	\caption{How to Import the package}
	
\end{code}

\subsection{Object creation}

\begin{itemize}
	\item A pandas Series \texttt{s} is created with some values, including a NaN (not a number) value.
	\item A sequence of dates is generated using the \texttt{pd.date\_range()} function and stored in the \texttt{dates} variable.
	\item A DataFrame \texttt{df} is created with random values using the \texttt{np.random.randn()} function. The DataFrame has 6 rows and 4 columns, with the dates as the index and column labels as 'A', 'B', 'C', and 'D'.
\end{itemize}

\begin{code}[h!]
	\lstinputlisting[language=Python, numbers=none, linerange={41-48}]{Code/Pandas/Pandas.py}    
	
	\caption{Object Creation Example}
	
\end{code}

\begin{verbatim}
	0    1.0
	1    3.0
	2    5.0
	3    NaN
	4    6.0
	5    8.0
	dtype: float64
	DatetimeIndex(['2014-01-01', '2014-01-02', '2014-01-03', '2014-01-04',
	'2014-01-05', '2014-01-06'],
	dtype='datetime64[ns]', freq='D')
	A         B         C         D
	2014-01-01 -1.473364  1.116274 -1.815663  0.198003
	2014-01-02 -0.094806 -0.501628 -1.722026  0.190776
	2014-01-03  0.163170  0.995486  0.854393 -2.657105
	2014-01-04  0.064040 -0.041879  0.581580 -0.896344
	2014-01-05 -0.209241 -0.910417  0.517225  0.097267
	2014-01-06  0.669742  0.226435 -0.609790 -0.741223
\end{verbatim}

\subsection{Viewing data}

\begin{itemize}
	
	\item The \texttt{head()} method is used to display the first few rows of the DataFrame \texttt{df}.
	
	\item The \texttt{tail()} method is used to display the last 3 rows of \texttt{df}.
	
	\item The \texttt{describe()} method provides a statistical summary of the DataFrame.
	
	\item The \texttt{sort\_values()} method is used to sort the DataFrame \texttt{df} by the values in column 'B'.
	
\end{itemize}

\begin{code}[h!]
	\lstinputlisting[language=Python, numbers=none, linerange={50-52}]{Code/Pandas/Pandas.py}    
	
	\caption{Viewing Data}
	
\end{code}

The results are not shown from now on due to the shortage of space. The manual for the user is provided, so by runing that, the results would appear in the console.

\subsection{Selection}

\subsubsection{Getting}

\begin{itemize}
	
	\item Column 'A' of the DataFrame \texttt{df} is selected using \texttt{df["A"]}.
	
	\item The first three rows of the DataFrame \texttt{df} are selected using slicing with \texttt{df[0:3]}.
	
\end{itemize}

\subsubsection{Selection by label}
\begin{itemize}
	
	\item Using label-based indexing with \texttt{df.loc}, all rows and columns 'A' and 'B' are selected.
	
\end{itemize}

\subsubsection{Selection by position}
\begin{itemize}
	\item The fourth row of the DataFrame \texttt{df} is selected using \texttt{df.iloc[3]}.
	\item Rows 3 to 4 and columns 0 to 1 are selected using \texttt{df.iloc[3:5, 0:2]}.
\end{itemize}

\subsubsection{Boolean indexing}
\begin{itemize}
	\item The DataFrame \texttt{df} is filtered using a Boolean condition \texttt{df["A"] > 0}.
\end{itemize}

\begin{code}[h!]
	\lstinputlisting[language=Python, numbers=none, linerange={54-69}]{Code/Pandas/Pandas.py}    
	
	\caption{Boolean Indexing Example}
	
\end{code}

\subsection{Setting}
\begin{itemize}
	\item The value at row 0, column 1 of the DataFrame \texttt{df} is modified to 0 using \texttt{df.iloc[0, 1] = 0}.
\end{itemize}

\begin{code}[h!]
	\lstinputlisting[language=Python, numbers=none, linerange={71-73}]{Code/Pandas/Pandas.py}    
	
	\caption{Setting Example}
	
\end{code}


\subsection{Handling missing data}
\begin{itemize}
	
	\item A new DataFrame \texttt{df1} is created by reindexing \texttt{df} and adding a new column 'E'.
	
	\item Values in column 'E' for the first two rows of \texttt{df1} are set to 1.
	
	\item The \texttt{dropna()} method is used to drop rows with any NaN values from \texttt{df1}.
	
	\item The \texttt{pd.isna()} function is used to check for NaN values in \texttt{df1}.
	
\end{itemize}


\begin{code}[h!]
	\lstinputlisting[language=Python, numbers=none, linerange={75-77}]{Code/Pandas/Pandas.py}    
	
	\caption{Handling Missing Data}
	
\end{code}


\subsection{Operations}

\begin{itemize}
	\item The \texttt{apply()} method is used to calculate the difference between the maximum and minimum values of each column in \texttt{df}.
\end{itemize}

\begin{code}[h!]
	\lstinputlisting[language=Python, numbers=none, linerange={79-80}]{Code/Pandas/Pandas.py}    
	
	\caption{Operations Example}
	
\end{code}

\subsection{Merge}

\begin{itemize}
	\item The \texttt{concat()} function is used to concatenate three chunks of the DataFrame \texttt{df}.
	
	\item Two DataFrames, \texttt{left} and \texttt{right}, are created with a common column 'key'.
	
	\item The \texttt{merge()} function is used to merge \texttt{left} and \texttt{right} DataFrames based on the 'key' column.
	
\end{itemize}

\begin{code}[h!]
	\lstinputlisting[language=Python, numbers=none, linerange={83-97}]{Code/Pandas/Pandas.py}    
	
	\caption{Merge Example}
	
\end{code}


\subsection{Grouping}

\begin{itemize}
	\item A DataFrame \texttt{df} is created with 'Animal' and 'Max Speed' columns.
	
	\item The \texttt{groupby()} method is used to group the DataFrame \texttt{df} by the 'Animal' column and calculate the mean of 'Max Speed' for each group.
	
\end{itemize}

\begin{code}[h!]
	\lstinputlisting[language=Python, numbers=none, linerange={99-104}]{Code/Pandas/Pandas.py}    
	
	\caption{Grouping Example}
	
\end{code}

\subsection{Reshaping}

\begin{itemize}
	
	\item A DataFrame \texttt{df\_single\_level\_cols} is created with two rows, two columns, and single-level column labels.
	
	\item The \texttt{stack()} method is used to stack the DataFrame, returning a Series.
	
	\item The \texttt{unstack()} method is used to unstack the Series, reverting it.
	
\end{itemize}

\begin{code}[h!]
	\lstinputlisting[language=Python, numbers=none, linerange={106-115}]{Code/Pandas/Pandas.py}    
	
	\caption{Reshaping Example}
	
\end{code}

\subsection{Importing and exporting data}

\begin{enumerate}
	\item \texttt{df.to\_csv("foo.csv")}: This line takes a Pandas DataFrame, \texttt{df}, and exports it to a CSV file named "foo.csv". The \texttt{to\_csv} method is used to convert the DataFrame into a CSV format and save it as a file. The resulting CSV file will contain the data from the DataFrame, with each row representing a data entry and each column representing a variable.
	
	\item \texttt{pd.read\_csv("foo.csv")}: This line reads the contents of the CSV file "foo.csv" and creates a new DataFrame using the data from the file. The \texttt{read\_csv} function from the Pandas library is used to read the CSV file and convert it back into a DataFrame.
\end{enumerate}

\begin{code}[h!]
	\lstinputlisting[language=Python, numbers=none, linerange={117-119}]{Code/Pandas/Pandas.py}    
	
	\caption{Reading CSV file }
	
\end{code}

\subsection{Pandas Error Handling}

Pandas provides various error handling mechanisms to handle exceptions and errors that may occur during data manipulation and analysis. These error handling techniques help in diagnosing and addressing issues that arise when working with data in Pandas. Some common error handling techniques in Pandas include:

\begin{enumerate}
	\item \textbf{Try-Except Blocks:} You can use standard Python try-except blocks to catch and handle specific exceptions that may occur during Pandas operations. For example, you can wrap a Pandas function call within a try block and use except blocks to handle specific exceptions, such as \texttt{ValueError} or \texttt{KeyError}, and perform appropriate error handling actions.
	
	\begin{code}[h!]
		\lstinputlisting[language=Python, numbers=none, linerange={121-130}]{Code/Pandas/Pandas.py}    
		
		\caption{Try Except Blocks}
		
	\end{code}
	
	\item \textbf{Error Reporting:} Pandas provides descriptive error messages that provide information about the nature of the error and the location where it occurred. These error messages can be useful for diagnosing and debugging issues in the data or the code. When an error occurs, Pandas displays an error message along with a traceback that helps identify the source of the error.
	
	\item \textbf{Error Handling Functions:} Pandas provides specific functions to handle errors and exceptions. For example, the \texttt{pd.options.mode} function allows you to set error handling modes, such as \texttt{raise} to raise exceptions for errors, \texttt{warn} to display warning messages instead of raising exceptions, and \texttt{ignore} to ignore errors and proceed with the operation.
	
	\begin{code}[h!]
		\lstinputlisting[language=Python, numbers=none, linerange={132-140}]{Code/Pandas/Pandas.py}    
		
		\caption{Error Handling Functions}
		
	\end{code}
	
	\item \textbf{Error Handling with DataFrame and Series:} Pandas provides methods and properties for error handling specific to DataFrame and Series objects. For example, the \texttt{.at} and \texttt{.iat} attributes allow for safe access and assignment of scalar values, raising exceptions if the provided index is not found. The \texttt{.get()} method allows retrieving values with a default value, avoiding exceptions when the key is not found.
	
	\begin{code}[h!]
		\lstinputlisting[language=Python, numbers=none, linerange={142-150}]{Code/Pandas/Pandas.py}    
		
		\caption{Error Handling with DataFrame and Series}
		
	\end{code}
	
\end{enumerate}

By leveraging these error handling techniques, you can effectively handle exceptions and errors that may arise during data operations in Pandas, ensuring smooth data processing and analysis.


\section{Further Reading}

\subsection{A Comprehensive Overview of Pandas}

Explore pandas, a Python library tailored for structured data analysis in statistics, finance, and social sciences. This paper covers its design, features, and its role as a foundational layer in Python's statistical computing landscape. Discover how pandas complements the scientific Python stack and improves upon data manipulation tools from languages like R. Additionally, insights into future development and growth opportunities for statistical applications in Python are discussed \cite{McKinney:2011}. Find it here : \url{https://www.dlr.de/sc/portaldata/15/resources/dokumente/pyhpc2011/submissions/pyhpc2011_submission_9.pdf}

\subsection{pandas: a python data analysis library}

This book, written by Wes McKinney (the creator of Pandas) \cite{McKinney:2015}, is a valuable resource for learning Pandas. It covers various aspects of data manipulation, analysis, and visualization using Pandas. The book also explores practical examples and real-world use cases. 
Find it here: \ul{https://www.oreilly.com/library/view/python-for-data/9781491957653/}

\subsection{Pandas Library}

This chapter presents pandas, a powerful Python library that introduces the DataFrame data structure. Derived from NumPy arrays, pandas revolutionizes data analysis by offering a structured framework. The core elements, Series and DataFrame, enable the organization of varied data types into a unified structure, facilitating easy application of methods or functions to the entire dataset or specific segments. This capability expands the horizons of efficient and comprehensive data analysis \cite{Betancourt:2019}. 
Find it here: \url{https://link.springer.com/chapter/10.1007/978-1-4842-5001-3_3}

\subsection{Python for Data Analysis}

The provided Books page with the book titled "Python for Data Analysis" by Wes McKinney. This book extensively covers the pandas library, a key tool for data analysis in Python. The link allows access to the book for further exploration of pandas and its applications in the context of data analysis using Python. \cite{McKinney:2012}. 
Access it here: \url{https://books.google.de/books?hl=en&lr=&id=v3n4_AK8vu0C&oi=fnd&pg=PR3&dq=pandas+library&ots=riDL2lxAsD&sig=652ON3D9JVUQ4BBtgvBuo-4tdgc&redir_esc=y#v=onepage&q=pandas%20library&f=false}

\subsection{Pandas Learning}
The provided link directs to a Google Books page for the book titled "Python for Data Science For Dummies" by John Paul Mueller and Luca Massaron. This book likely covers various aspects of Python for data science, including the pandas library. The link allows access to the book for further exploration of pandas and its applications within the broader context of data science using Python \cite{Heydt:2017}. 
You can find the official Pandas learning book here: \url{https://books.google.de/books?hl=en&lr=&id=Sng5DwAAQBAJ&oi=fnd&pg=PP1&dq=pandas+library&ots=J-L7JhMo75&sig=9RaIsH-obmG0vNc9_qjL6AnNUHA&redir_esc=y#v=onepage&q=pandas%20library&f=false}

Start by exploring the official documentation, as it provides the most current and comprehensive information on Pandas. The books, online resources, and video tutorials mentioned above will further enhance your understanding and offer practical insights into using Pandas effectively. Additionally, regularly experimenting with real-world data and engaging with the Pandas community through forums or blogs can help solidify your knowledge and keep you updated on new features and best practices.




